% part: intuitionistic-logic
% chapter: propositions-as-types

\documentclass[../../../include/open-logic-chapter]{subfiles}

\begin{document}

\olchapter{int}{pty}{قضیے بطور نوعاں}

\begin{editorial}
  ایہ کری–ہاورڈ میل اُتے اک باب دا \emph{بہت تجرباتی} مسودہ اے۔ ایہنوں ہور سمجھاؤنی تے وجہ دسّن دی لوڑ اے، تے شاید ایہدے وچ غلطیاں تے رہندیاں گلاں ہون۔ نارمل روپ وچ لیاؤن دے ثبوت نوں مُڑ ویکھنا تے ودھاؤنا چاہیدا اے۔ حاصلِ ضرب دی نوع لئی کوئی مثال نہیں۔ تھاں بدلّن والیاں تے سادہ کرن والیاں تبدیلیاں شامل نہیں۔ جدوں (نوع بندی والے) لامبڈا کیلکولس بارے مواد وی ہووے گا، جیہڑا ایتھے بنیادی طور اُتے پہلاں توں منّیا گیا اے، تاں ایہ ہور سوکھا سمجھ آوے گا۔ بہت احتیاط نال ورتو۔
\end{editorial}

\olimport{introduction}
%\olimport{sequent-natural-deduction}
\olimport{proof-terms}
\olimport{proofs-to-terms}
\olimport{terms-to-proofs}
\olimport{types}
\olimport{reduction}
\olimport{type-preservation}
\olimport{normalization}

\OLEndChapterHook

\end{document}
